\section{Calibration with a Synthetic Examinee Population}
\label{sec:calibration}

% The 15-profile logit-b path (legacy bank pipeline), plus the bridge to
% the 45-cell Phase 1 grid used for the frozen bank.
With no human test-takers available, difficulty has to be estimated from
the responses of a \emph{synthetic} examinee population. I use two
related estimators. The first is a fast pass-rate logit used during item
generation. The second is the full Rasch fit that produces the frozen
bank, which I describe in Section~\ref{sec:pipeline}.

\subsection{Fifteen solver profiles and the pass-rate logit}
During generation, each candidate item is administered to 15 synthetic
solver profiles. The profiles span five expertise tiers (layperson,
life-sciences professional, regulatory-affairs coordinator,
regulatory-affairs specialist, and senior FDA auditor) at three sampling
temperatures each. Let $k_i$ be the number of profiles that pass item $i$
and let $\hat{P}_i = k_i / 15$ be the empirical pass rate. Under a Rasch
model in which the profile population is centered at $\theta = 0$, the
probability of a correct response to an item of difficulty $b$ is
$\sigma(\theta - b)$. So at $\theta = 0$ we have
$\hat{P}_i = \sigma(-b_i)$, and therefore

\begin{equation}
\hat{b}_i \;=\; \ln\!\left(\frac{1 - \hat{P}_i}{\hat{P}_i}\right),
\qquad \hat{P}_i \in [0.01,\, 0.99],
\label{eq:logitb}
\end{equation}

where $\hat{P}_i$ is clipped to keep $\hat{b}_i$ finite at the extremes.
Equation~\eqref{eq:logitb} ranges from $b \approx +4.6$ for a very hard
item ($\hat{P}=0.01$) to $b \approx -4.6$ for a trivially easy one
($\hat{P}=0.99$).

\subsection{Retention thresholds}
Items are retained or discarded by pass rate, as shown in
Table~\ref{tab:retention}. I widened the thresholds to $[0.05, 0.95]$ in
order to capture items that are near the extremes but still informative.

\begin{table}[t]
\centering
\small
\begin{tabular}{@{}llc@{}}
\toprule
pass rate $\hat{P}$ & outcome & retained? \\
\midrule
$\hat{P} < 0.05$        & \texttt{discarded\_too\_hard} & no \\
$0.05 \le \hat{P} \le 0.95$ & \texttt{retained}        & yes \\
$\hat{P} > 0.95$        & \texttt{retained\_easy}       & yes \\
\bottomrule
\end{tabular}
\caption{Pass-rate retention rule. Difficulty $b$ from
Eq.~\eqref{eq:logitb} is recorded for all items, including the
too-hard ones, so that they can still enter ability estimation as
failures.}
\label{tab:retention}
\end{table}

\subsection{Key assumption and its risk}
Equation~\eqref{eq:logitb} assumes that the synthetic population is
centered at $\theta = 0$. This assumption carries a risk. If the profiles
lean expert on average ($\bar{\theta} > 0$), then $b$ is biased downward
and the items look easier than they really are. Also, the 15 profiles are
highly correlated. They are the same base model run under different
prompts and temperatures, so they are not 15 independent test-takers. The
full Rasch fit in Section~\ref{sec:pipeline} relaxes the
centered-population assumption, because it estimates examinee abilities
jointly and anchors them to a fixed baseline. The remaining independence
problem is discussed in Section~\ref{sec:discussion}.
